\chapter{Literature Review}
\label{ch:litreview}

\section{Protein--protein affinity prediction}

Protein--protein affinity prediction remains less standardized than several neighbouring tasks in structural bioinformatics \citep{Guo2022PPBAReview,Liu2024PPBAffinity}. Much of the older supervised literature focused on mutation-induced changes in binding energy, supported by resources such as SKEMPI 2.0, rather than on direct prediction of absolute binding affinity for heterogeneous wild-type complexes \citep{Jankauskaite2019SKEMPI}. Absolute-affinity benchmarks do exist, including the structure-based benchmark assembled by Kastritis and coworkers, but they are small by current deep-learning standards \citep{Kastritis2011Benchmark}. PPB-Affinity is therefore an important recent development because it was assembled specifically for direct protein--protein affinity prediction at larger scale \citep{Liu2024PPBAffinity}.

This task is difficult for straightforward physical reasons. The target is thermodynamic, but the input is usually a single experimentally observed structure. Reviews of the field emphasize that limited labelled data, structural heterogeneity, and the gap between one structure and an underlying binding ensemble remain major obstacles \citep{Guo2022PPBAReview,Liu2024PPBAffinity}.

\section{Physics-based and empirical approaches}

One response is to estimate affinity directly from simulation. Umbrella sampling and alchemical free-energy methods are established routes to free-energy estimation, but they are computationally expensive and sensitive to protocol design \citep{Torrie1977Umbrella,Cournia2017RBFE}. At lower cost, MM/PBSA and MM/GBSA provide widely used endpoint approximations, although their performance depends strongly on solvent treatment, entropy approximations, and state averaging \citep{Genheden2015MMGBSA}. This literature motivates taking molecular dynamics seriously, but it does not make full free-energy calculation the only practical option.

Another response is to predict affinity directly from structure. PRODIGY is a representative example of a fast empirical model built from interfacial contact descriptors \citep{Vangone2015Contacts}. Other earlier approaches similarly relied on engineered structural features and conventional regressors \citep{Moal2011AffinityPrediction}. More recent work has moved toward learned structural representations. The PPB-Affinity paper includes a benchmark deep model on the new dataset \citep{Liu2024PPBAffinity}, and recent studies such as PPI-Graphomer have explored pretrained and graph-based models for direct affinity prediction \citep{Xie2025PPIGraphomer}.

\section{Why the protein--ligand comparison matters}

The protein--ligand literature provides a useful contrast. Datasets such as PDBbind have given that field a longer and more standardized benchmark tradition \citep{Wang2004PDBbind}. This has supported a broad range of learned structural models, including voxelized approaches such as Pafnucy \citep{Stepniewska2018Pafnucy} and more recent trajectory-aware models trained on MD-derived data \citep{Libouban2025Spatiotemporal}. The protein--ligand literature therefore provides concrete examples of successful geometric and trajectory-aware affinity models.

Protein--protein complexes make the representation question less obvious. A ligand naturally defines a local binding object inside a pocket; a protein--protein complex does not. Even before model choice, one must decide whether the relevant representation is the whole complex, the interface only, or some local patch around the interface. That is a central design question in this thesis.

\section{Why use a variational graph model}

Variational autoencoders are well established in molecular representation learning as a way to combine reconstruction with latent-space regularization \citep{Kingma2014VAE,GomezBombarelli2018MolecularVAE,Kipf2016VGAE,Simonovsky2018GraphVAE}. They are not the default standard in direct protein--protein affinity prediction, where supervised regression remains more common \citep{Guo2022PPBAReview,Xie2025PPIGraphomer}. They are still relevant here because the thesis asks not only which predictor is strongest, but also whether a reconstruction-based latent representation can capture useful structure in residue graphs enriched with molecular-dynamics-derived features.

\section{Where this thesis fits}

The literature suggests a clear gap. Rigorous free-energy methods are principled but expensive. Fast empirical scorers are cheap but compressed. Learned structural models are becoming more capable, but direct protein--protein affinity prediction still lacks a single settled representation or modelling strategy. This thesis therefore asks a focused question: starting from a residue-graph representation of a protein complex, do short explicit-solvent molecular-dynamics summaries improve held-out affinity prediction, and does that answer depend on graph view or model family?
